% --- Archon physics formalization source begin ---
% archon:physics
% archon:covers PhyXMiniProblems/problem_phyx_mini_0894.lean
% archon:source-report reports/phyx_mini/problem_phyx_mini_0894.source.json

\chapter{Physics problem phyx\_mini\_0894}
\label{ch:PhyXMiniProblems_problem_phyx_mini_0894}

\paragraph{Problem source.}
\#\# Question

What is the magnitude of the force \textbackslash{}( \textbackslash{}vec\{F\} \textbackslash{}) on the \textbackslash{}( 1.0\textbackslash{},\textbackslash{}mathrm\{nC\} \textbackslash{}) charge in figure?

\#\# Answer choices

- A: A: 1.25 \textbackslash{}times 10\textasciicircum{}\{-3\}N

- B: B: 1.35 \textbackslash{}times 10\textasciicircum{}\{-3\}N

- C: C: 1.8 \textbackslash{}times 10\textasciicircum{}\{-4\}N

- D: D: 1.33 \textbackslash{}times 10\textasciicircum{}\{-3\}N

\#\# Figure caption (auxiliary; use the image as primary evidence)

The image is a diagram of an equilateral triangle with three charged particles at its vertices. 

1. **Charges:**
   - The top vertex has a positive charge labeled \textbackslash{}(1.0 \textbackslash{}, \textbackslash{}text\{nC\}\textbackslash{}).
   - The bottom left vertex has a positive charge labeled \textbackslash{}(2.0 \textbackslash{}, \textbackslash{}text\{nC\}\textbackslash{}).
   - The bottom right vertex has a negative charge labeled \textbackslash{}(-2.0 \textbackslash{}, \textbackslash{}text\{nC\}\textbackslash{}).

2. **Distances:**
   - Each side of the triangle measures \textbackslash{}(1.0 \textbackslash{}, \textbackslash{}text\{cm\}\textbackslash{}).

3. **Angles:**
   - The angles between each side are \textbackslash{}(60\textasciicircum{}\textbackslash{}circ\textbackslash{}), indicating it is an equilateral triangle. 

4. **Connections:**
   - Dotted lines connect the vertices, emphasizing the shape and distance between the charges.

\#\# Dataset metadata

Electrostatics; Spatial Relation Reasoning, Multi-Formula Reasoning

\paragraph{Recorded answer/context.}
C: C: 1.8 \textbackslash{}times 10\textasciicircum{}\{-4\}N

\paragraph{Figure/image path.}
/root/proposal\_for\_physic/science-mango/phyx\_mini\_run/phyx\_data/test\_image/894.png

\paragraph{Formalization target.}
create a compiling Lean file with sorry bodies at `PhyXMiniProblems/problem\_phyx\_mini\_0894.lean`. The Lean declarations must preserve the physical quantities, dimensions or dimensional roles, figure labels, governing-law hypotheses, and final relation expressed by this problem.
Use Mathlib/Physlib names found through LeanExplore where available. If a physics API is missing, introduce faithful local abstractions rather than scalar placeholder aliases.

\begin{theorem}[Physics formalization target]
\label{thm:physics:phyx_mini_0894:target}
\lean{PhyXMiniProblems.ProblemPhyXMini0894.problem_phyx_mini_0894}
\uses{def:physics:phyx-mini-0894:phyxminiproblems-problemphyxmini0894-forcemagnitudeinnewtons, def:physics:phyx-mini-0894:phyxminiproblems-problemphyxmini0894-equilateralpointchargesetup, def:physics:phyx-mini-0894:phyxminiproblems-problemphyxmini0894-matchessuppliedequilateralchargefigure, def:physics:phyx-mini-0894:phyxminiproblems-problemphyxmini0894-hasequilateraltrianglegeometry, def:physics:phyx-mini-0894:phyxminiproblems-problemphyxmini0894-hasphysicalpointchargeparameters, def:physics:phyx-mini-0894:phyxminiproblems-problemphyxmini0894-usesroundedcoulombconstant, def:physics:phyx-mini-0894:phyxminiproblems-problemphyxmini0894-satisfiespointchargecoulombforcelaw, def:physics:phyx-mini-0894:phyxminiproblems-problemphyxmini0894-satisfieselectricforcesuperposition, def:physics:phyx-mini-0894:phyxminiproblems-problemphyxmini0894-recordeddatasetanswer}
The assigned autoformalize agent should translate this physics problem into Lean declarations in the covered file, with theorem and lemma proof bodies written as `by sorry`.
\end{theorem}
\begin{proof}
This is an autoformalization task, not a proof task. Produce statements that can later be proved by the physics prover without weakening the physical model.
\end{proof}

% --- Archon declaration topology begin ---
\section{Declaration topology}
\begin{definition}[Planar Vector]
  \label{def:physics:phyx-mini-0894:phyxminiproblems-problemphyxmini0894-planarvector}
  \lean{PhyXMiniProblems.ProblemPhyXMini0894.PlanarVector}
  A planar scalar readout, used for metre coordinates and force components
\end{definition}
\begin{proof}
  This is the declared model definition of Planar Vector.
\end{proof}
\begin{definition}[force Dimension]
  \label{def:physics:phyx-mini-0894:phyxminiproblems-problemphyxmini0894-forcedimension}
  \lean{PhyXMiniProblems.ProblemPhyXMini0894.forceDimension}
  The physical dimension of force, M L T\textasciicircum{}-2
\end{definition}
\begin{proof}
  This is the declared model definition of force Dimension.
\end{proof}
\begin{definition}[Length Quantity]
  \label{def:physics:phyx-mini-0894:phyxminiproblems-problemphyxmini0894-lengthquantity}
  \lean{PhyXMiniProblems.ProblemPhyXMini0894.LengthQuantity}
  A nonnegative, unit-independent physical length
\end{definition}
\begin{proof}
  This is the declared model definition of Length Quantity.
\end{proof}
\begin{definition}[Signed Charge Quantity]
  \label{def:physics:phyx-mini-0894:phyxminiproblems-problemphyxmini0894-signedchargequantity}
  \lean{PhyXMiniProblems.ProblemPhyXMini0894.SignedChargeQuantity}
  A signed, unit-independent physical electric charge
\end{definition}
\begin{proof}
  This is the declared model definition of Signed Charge Quantity.
\end{proof}
\begin{definition}[Planar Position Quantity]
  \label{def:physics:phyx-mini-0894:phyxminiproblems-problemphyxmini0894-planarpositionquantity}
  \lean{PhyXMiniProblems.ProblemPhyXMini0894.PlanarPositionQuantity}
  \uses{def:physics:phyx-mini-0894:phyxminiproblems-problemphyxmini0894-planarvector}
  A unit-independent planar position carrying the dimension of length
\end{definition}
\begin{proof}
  This is the declared model definition of Planar Position Quantity.
\end{proof}
\begin{definition}[Planar Force Quantity]
  \label{def:physics:phyx-mini-0894:phyxminiproblems-problemphyxmini0894-planarforcequantity}
  \lean{PhyXMiniProblems.ProblemPhyXMini0894.PlanarForceQuantity}
  \uses{def:physics:phyx-mini-0894:phyxminiproblems-problemphyxmini0894-planarvector, def:physics:phyx-mini-0894:phyxminiproblems-problemphyxmini0894-forcedimension}
  A unit-independent planar physical force vector
\end{definition}
\begin{proof}
  This is the declared model definition of Planar Force Quantity.
\end{proof}
\begin{definition}[length In Meters]
  \label{def:physics:phyx-mini-0894:phyxminiproblems-problemphyxmini0894-lengthinmeters}
  \lean{PhyXMiniProblems.ProblemPhyXMini0894.lengthInMeters}
  \uses{def:physics:phyx-mini-0894:phyxminiproblems-problemphyxmini0894-lengthquantity}
  Coherent-SI metre readout of a physical length
\end{definition}
\begin{proof}
  This is the declared model definition of length In Meters.
\end{proof}
\begin{definition}[length In Centimeters]
  \label{def:physics:phyx-mini-0894:phyxminiproblems-problemphyxmini0894-lengthincentimeters}
  \lean{PhyXMiniProblems.ProblemPhyXMini0894.lengthInCentimeters}
  \uses{def:physics:phyx-mini-0894:phyxminiproblems-problemphyxmini0894-lengthquantity, def:physics:phyx-mini-0894:phyxminiproblems-problemphyxmini0894-lengthinmeters}
  Centimetre readout used by the three printed side labels
\end{definition}
\begin{proof}
  This is the declared model definition of length In Centimeters.
\end{proof}
\begin{definition}[charge In Coulombs]
  \label{def:physics:phyx-mini-0894:phyxminiproblems-problemphyxmini0894-chargeincoulombs}
  \lean{PhyXMiniProblems.ProblemPhyXMini0894.chargeInCoulombs}
  \uses{def:physics:phyx-mini-0894:phyxminiproblems-problemphyxmini0894-signedchargequantity}
  Coherent-SI coulomb readout of a signed electric charge
\end{definition}
\begin{proof}
  This is the declared model definition of charge In Coulombs.
\end{proof}
\begin{definition}[charge In Nanocoulombs]
  \label{def:physics:phyx-mini-0894:phyxminiproblems-problemphyxmini0894-chargeinnanocoulombs}
  \lean{PhyXMiniProblems.ProblemPhyXMini0894.chargeInNanocoulombs}
  \uses{def:physics:phyx-mini-0894:phyxminiproblems-problemphyxmini0894-signedchargequantity, def:physics:phyx-mini-0894:phyxminiproblems-problemphyxmini0894-chargeincoulombs}
  Nanocoulomb readout used by the three printed charge labels
\end{definition}
\begin{proof}
  This is the declared model definition of charge In Nanocoulombs.
\end{proof}
\begin{definition}[position In Meters]
  \label{def:physics:phyx-mini-0894:phyxminiproblems-problemphyxmini0894-positioninmeters}
  \lean{PhyXMiniProblems.ProblemPhyXMini0894.positionInMeters}
  \uses{def:physics:phyx-mini-0894:phyxminiproblems-problemphyxmini0894-planarvector, def:physics:phyx-mini-0894:phyxminiproblems-problemphyxmini0894-planarpositionquantity}
  Coherent-SI planar metre readout of a physical position
\end{definition}
\begin{proof}
  This is the declared model definition of position In Meters.
\end{proof}
\begin{definition}[force Vector In Newtons]
  \label{def:physics:phyx-mini-0894:phyxminiproblems-problemphyxmini0894-forcevectorinnewtons}
  \lean{PhyXMiniProblems.ProblemPhyXMini0894.forceVectorInNewtons}
  \uses{def:physics:phyx-mini-0894:phyxminiproblems-problemphyxmini0894-planarvector, def:physics:phyx-mini-0894:phyxminiproblems-problemphyxmini0894-planarforcequantity}
  Coherent-SI planar newton readout of a physical force vector
\end{definition}
\begin{proof}
  This is the declared model definition of force Vector In Newtons.
\end{proof}
\begin{definition}[force Magnitude In Newtons]
  \label{def:physics:phyx-mini-0894:phyxminiproblems-problemphyxmini0894-forcemagnitudeinnewtons}
  \lean{PhyXMiniProblems.ProblemPhyXMini0894.forceMagnitudeInNewtons}
  \uses{def:physics:phyx-mini-0894:phyxminiproblems-problemphyxmini0894-planarforcequantity, def:physics:phyx-mini-0894:phyxminiproblems-problemphyxmini0894-forcevectorinnewtons}
  Magnitude in newtons of a dimensionful planar force vector
\end{definition}
\begin{proof}
  This is the declared model definition of force Magnitude In Newtons.
\end{proof}
\begin{definition}[Triangle Vertex]
  \label{def:physics:phyx-mini-0894:phyxminiproblems-problemphyxmini0894-trianglevertex}
  \lean{PhyXMiniProblems.ProblemPhyXMini0894.TriangleVertex}
  The three charge sites in image 894.png
\end{definition}
\begin{proof}
  This is the declared model definition of Triangle Vertex.
\end{proof}
\begin{definition}[Lower Source]
  \label{def:physics:phyx-mini-0894:phyxminiproblems-problemphyxmini0894-lowersource}
  \lean{PhyXMiniProblems.ProblemPhyXMini0894.LowerSource}
  The two lower charges that exert force on the top charge
\end{definition}
\begin{proof}
  This is the declared model definition of Lower Source.
\end{proof}
\begin{definition}[vertex]
  \label{def:physics:phyx-mini-0894:phyxminiproblems-problemphyxmini0894-lowersource-vertex}
  \lean{PhyXMiniProblems.ProblemPhyXMini0894.LowerSource.vertex}
  \uses{def:physics:phyx-mini-0894:phyxminiproblems-problemphyxmini0894-trianglevertex, def:physics:phyx-mini-0894:phyxminiproblems-problemphyxmini0894-lowersource}
  Vertex occupied by each source charge acting on the top charge
\end{definition}
\begin{proof}
  This is the declared model definition of vertex.
\end{proof}
\begin{definition}[Triangle Edge]
  \label{def:physics:phyx-mini-0894:phyxminiproblems-problemphyxmini0894-triangleedge}
  \lean{PhyXMiniProblems.ProblemPhyXMini0894.TriangleEdge}
  The three dashed sides of the triangular diagram
\end{definition}
\begin{proof}
  This is the declared model definition of Triangle Edge.
\end{proof}
\begin{definition}[endpoints]
  \label{def:physics:phyx-mini-0894:phyxminiproblems-problemphyxmini0894-triangleedge-endpoints}
  \lean{PhyXMiniProblems.ProblemPhyXMini0894.TriangleEdge.endpoints}
  \uses{def:physics:phyx-mini-0894:phyxminiproblems-problemphyxmini0894-trianglevertex, def:physics:phyx-mini-0894:phyxminiproblems-problemphyxmini0894-triangleedge}
  Ordered endpoints of each triangle side
\end{definition}
\begin{proof}
  This is the declared model definition of endpoints.
\end{proof}
\begin{definition}[Printed Angle Site]
  \label{def:physics:phyx-mini-0894:phyxminiproblems-problemphyxmini0894-printedanglesite}
  \lean{PhyXMiniProblems.ProblemPhyXMini0894.PrintedAngleSite}
  The two lower vertices carrying explicit 60 degree arcs in the bitmap
\end{definition}
\begin{proof}
  This is the declared model definition of Printed Angle Site.
\end{proof}
\begin{definition}[Charge Circle Color]
  \label{def:physics:phyx-mini-0894:phyxminiproblems-problemphyxmini0894-chargecirclecolor}
  \lean{PhyXMiniProblems.ProblemPhyXMini0894.ChargeCircleColor}
  Fill colours used by the charge circles in the primary image
\end{definition}
\begin{proof}
  This is the declared model definition of Charge Circle Color.
\end{proof}
\begin{definition}[Charge Sign Glyph]
  \label{def:physics:phyx-mini-0894:phyxminiproblems-problemphyxmini0894-chargesignglyph}
  \lean{PhyXMiniProblems.ProblemPhyXMini0894.ChargeSignGlyph}
  Sign glyph drawn inside a charge circle
\end{definition}
\begin{proof}
  This is the declared model definition of Charge Sign Glyph.
\end{proof}
\begin{definition}[expected Charge In Nanocoulombs]
  \label{def:physics:phyx-mini-0894:phyxminiproblems-problemphyxmini0894-expectedchargeinnanocoulombs}
  \lean{PhyXMiniProblems.ProblemPhyXMini0894.expectedChargeInNanocoulombs}
  \uses{def:physics:phyx-mini-0894:phyxminiproblems-problemphyxmini0894-trianglevertex}
  Expected signed nanocoulomb label at each figure vertex
\end{definition}
\begin{proof}
  This is the declared model definition of expected Charge In Nanocoulombs.
\end{proof}
\begin{definition}[expected Charge Color]
  \label{def:physics:phyx-mini-0894:phyxminiproblems-problemphyxmini0894-expectedchargecolor}
  \lean{PhyXMiniProblems.ProblemPhyXMini0894.expectedChargeColor}
  \uses{def:physics:phyx-mini-0894:phyxminiproblems-problemphyxmini0894-trianglevertex, def:physics:phyx-mini-0894:phyxminiproblems-problemphyxmini0894-chargecirclecolor}
  Expected charge-circle colour at each figure vertex
\end{definition}
\begin{proof}
  This is the declared model definition of expected Charge Color.
\end{proof}
\begin{definition}[expected Charge Glyph]
  \label{def:physics:phyx-mini-0894:phyxminiproblems-problemphyxmini0894-expectedchargeglyph}
  \lean{PhyXMiniProblems.ProblemPhyXMini0894.expectedChargeGlyph}
  \uses{def:physics:phyx-mini-0894:phyxminiproblems-problemphyxmini0894-trianglevertex, def:physics:phyx-mini-0894:phyxminiproblems-problemphyxmini0894-chargesignglyph}
  Expected sign glyph at each figure vertex
\end{definition}
\begin{proof}
  This is the declared model definition of expected Charge Glyph.
\end{proof}
\begin{definition}[Equilateral Charge Figure]
  \label{def:physics:phyx-mini-0894:phyxminiproblems-problemphyxmini0894-equilateralchargefigure}
  \lean{PhyXMiniProblems.ProblemPhyXMini0894.EquilateralChargeFigure}
  \uses{def:physics:phyx-mini-0894:phyxminiproblems-problemphyxmini0894-trianglevertex, def:physics:phyx-mini-0894:phyxminiproblems-problemphyxmini0894-triangleedge, def:physics:phyx-mini-0894:phyxminiproblems-problemphyxmini0894-printedanglesite, def:physics:phyx-mini-0894:phyxminiproblems-problemphyxmini0894-chargecirclecolor, def:physics:phyx-mini-0894:phyxminiproblems-problemphyxmini0894-chargesignglyph}
  Literal presentation data transcribed from the supplied raster. It contains no force magnitude or answer-choice field
\end{definition}
\begin{proof}
  This is the declared model definition of Equilateral Charge Figure.
\end{proof}
\begin{definition}[Equilateral Point Charge Setup]
  \label{def:physics:phyx-mini-0894:phyxminiproblems-problemphyxmini0894-equilateralpointchargesetup}
  \lean{PhyXMiniProblems.ProblemPhyXMini0894.EquilateralPointChargeSetup}
  \uses{def:physics:phyx-mini-0894:phyxminiproblems-problemphyxmini0894-lengthquantity, def:physics:phyx-mini-0894:phyxminiproblems-problemphyxmini0894-signedchargequantity, def:physics:phyx-mini-0894:phyxminiproblems-problemphyxmini0894-planarpositionquantity, def:physics:phyx-mini-0894:phyxminiproblems-problemphyxmini0894-planarforcequantity, def:physics:phyx-mini-0894:phyxminiproblems-problemphyxmini0894-trianglevertex, def:physics:phyx-mini-0894:phyxminiproblems-problemphyxmini0894-lowersource, def:physics:phyx-mini-0894:phyxminiproblems-problemphyxmini0894-equilateralchargefigure}
  Independent physical quantities in the electrostatic setup. The net force and both pair forces are stored as independent dimensionful vectors; none is defined from a displayed answer
\end{definition}
\begin{proof}
  This is the declared model definition of Equilateral Point Charge Setup.
\end{proof}
\begin{definition}[displacement From Source To Top In Meters]
  \label{def:physics:phyx-mini-0894:phyxminiproblems-problemphyxmini0894-displacementfromsourcetotopinmeters}
  \lean{PhyXMiniProblems.ProblemPhyXMini0894.displacementFromSourceToTopInMeters}
  \uses{def:physics:phyx-mini-0894:phyxminiproblems-problemphyxmini0894-planarvector, def:physics:phyx-mini-0894:phyxminiproblems-problemphyxmini0894-positioninmeters, def:physics:phyx-mini-0894:phyxminiproblems-problemphyxmini0894-lowersource, def:physics:phyx-mini-0894:phyxminiproblems-problemphyxmini0894-equilateralpointchargesetup}
  Displacement in metres from a lower source to the top charge
\end{definition}
\begin{proof}
  This is the declared model definition of displacement From Source To Top In Meters.
\end{proof}
\begin{definition}[edge Length In Meters]
  \label{def:physics:phyx-mini-0894:phyxminiproblems-problemphyxmini0894-edgelengthinmeters}
  \lean{PhyXMiniProblems.ProblemPhyXMini0894.edgeLengthInMeters}
  \uses{def:physics:phyx-mini-0894:phyxminiproblems-problemphyxmini0894-positioninmeters, def:physics:phyx-mini-0894:phyxminiproblems-problemphyxmini0894-triangleedge, def:physics:phyx-mini-0894:phyxminiproblems-problemphyxmini0894-equilateralpointchargesetup}
  Metre distance between the endpoints of a named triangle edge
\end{definition}
\begin{proof}
  This is the declared model definition of edge Length In Meters.
\end{proof}
\begin{definition}[Matches Supplied Equilateral Charge Figure]
  \label{def:physics:phyx-mini-0894:phyxminiproblems-problemphyxmini0894-matchessuppliedequilateralchargefigure}
  \lean{PhyXMiniProblems.ProblemPhyXMini0894.MatchesSuppliedEquilateralChargeFigure}
  \uses{def:physics:phyx-mini-0894:phyxminiproblems-problemphyxmini0894-lengthincentimeters, def:physics:phyx-mini-0894:phyxminiproblems-problemphyxmini0894-chargeinnanocoulombs, def:physics:phyx-mini-0894:phyxminiproblems-problemphyxmini0894-expectedchargeinnanocoulombs, def:physics:phyx-mini-0894:phyxminiproblems-problemphyxmini0894-expectedchargecolor, def:physics:phyx-mini-0894:phyxminiproblems-problemphyxmini0894-expectedchargeglyph, def:physics:phyx-mini-0894:phyxminiproblems-problemphyxmini0894-equilateralpointchargesetup}
  Exact qualitative and numerical evidence from image 894.png, including the association of printed scalar labels with physical dimensionful quantities
\end{definition}
\begin{proof}
  This is the declared model definition of Matches Supplied Equilateral Charge Figure.
\end{proof}
\begin{definition}[Has Equilateral Triangle Geometry]
  \label{def:physics:phyx-mini-0894:phyxminiproblems-problemphyxmini0894-hasequilateraltrianglegeometry}
  \lean{PhyXMiniProblems.ProblemPhyXMini0894.HasEquilateralTriangleGeometry}
  \uses{def:physics:phyx-mini-0894:phyxminiproblems-problemphyxmini0894-lengthinmeters, def:physics:phyx-mini-0894:phyxminiproblems-problemphyxmini0894-positioninmeters, def:physics:phyx-mini-0894:phyxminiproblems-problemphyxmini0894-equilateralpointchargesetup, def:physics:phyx-mini-0894:phyxminiproblems-problemphyxmini0894-edgelengthinmeters}
  An oriented metre-coordinate realization of the equilateral triangle shown in the image: the base is horizontal, the top charge lies above its midpoint, and every edge has the common physical side length
\end{definition}
\begin{proof}
  This is the declared model definition of Has Equilateral Triangle Geometry.
\end{proof}
\begin{definition}[Has Physical Point Charge Parameters]
  \label{def:physics:phyx-mini-0894:phyxminiproblems-problemphyxmini0894-hasphysicalpointchargeparameters}
  \lean{PhyXMiniProblems.ProblemPhyXMini0894.HasPhysicalPointChargeParameters}
  \uses{def:physics:phyx-mini-0894:phyxminiproblems-problemphyxmini0894-lengthinmeters, def:physics:phyx-mini-0894:phyxminiproblems-problemphyxmini0894-equilateralpointchargesetup, def:physics:phyx-mini-0894:phyxminiproblems-problemphyxmini0894-displacementfromsourcetotopinmeters}
  Positivity and separation conditions selecting the physical branch
\end{definition}
\begin{proof}
  This is the declared model definition of Has Physical Point Charge Parameters.
\end{proof}
\begin{definition}[Uses Rounded Coulomb Constant]
  \label{def:physics:phyx-mini-0894:phyxminiproblems-problemphyxmini0894-usesroundedcoulombconstant}
  \lean{PhyXMiniProblems.ProblemPhyXMini0894.UsesRoundedCoulombConstant}
  \uses{def:physics:phyx-mini-0894:phyxminiproblems-problemphyxmini0894-equilateralpointchargesetup}
  The rounded constant 9.0 * 10\textasciicircum{}9 N m\textasciicircum{}2/C\textasciicircum{}2 used in the textbook multiple-choice calculation. This is an independent calibration, not a force answer
\end{definition}
\begin{proof}
  This is the declared model definition of Uses Rounded Coulomb Constant.
\end{proof}
\begin{definition}[Satisfies Point Charge Coulomb Force Law]
  \label{def:physics:phyx-mini-0894:phyxminiproblems-problemphyxmini0894-satisfiespointchargecoulombforcelaw}
  \lean{PhyXMiniProblems.ProblemPhyXMini0894.SatisfiesPointChargeCoulombForceLaw}
  \uses{def:physics:phyx-mini-0894:phyxminiproblems-problemphyxmini0894-chargeincoulombs, def:physics:phyx-mini-0894:phyxminiproblems-problemphyxmini0894-forcevectorinnewtons, def:physics:phyx-mini-0894:phyxminiproblems-problemphyxmini0894-equilateralpointchargesetup, def:physics:phyx-mini-0894:phyxminiproblems-problemphyxmini0894-displacementfromsourcetotopinmeters}
  For displacement r from a source to the top charge, Coulomb's vector law is F = k q\_top q\_source r / ‖r‖\textasciicircum{}3. Its signed charge product gives repulsion from the positive lower-left charge and attraction toward the negative lower-right charge. This premise contains no numerical net-force result
\end{definition}
\begin{proof}
  This is the declared model definition of Satisfies Point Charge Coulomb Force Law.
\end{proof}
\begin{definition}[Satisfies Electric Force Superposition]
  \label{def:physics:phyx-mini-0894:phyxminiproblems-problemphyxmini0894-satisfieselectricforcesuperposition}
  \lean{PhyXMiniProblems.ProblemPhyXMini0894.SatisfiesElectricForceSuperposition}
  \uses{def:physics:phyx-mini-0894:phyxminiproblems-problemphyxmini0894-forcevectorinnewtons, def:physics:phyx-mini-0894:phyxminiproblems-problemphyxmini0894-equilateralpointchargesetup}
  The force on the top charge is the vector sum of the two pair forces
\end{definition}
\begin{proof}
  This is the declared model definition of Satisfies Electric Force Superposition.
\end{proof}
\begin{definition}[Answer Choice]
  \label{def:physics:phyx-mini-0894:phyxminiproblems-problemphyxmini0894-answerchoice}
  \lean{PhyXMiniProblems.ProblemPhyXMini0894.AnswerChoice}
  The four force magnitudes printed in the answer list
\end{definition}
\begin{proof}
  This is the declared model definition of Answer Choice.
\end{proof}
\begin{definition}[force In Newtons]
  \label{def:physics:phyx-mini-0894:phyxminiproblems-problemphyxmini0894-answerchoice-forceinnewtons}
  \lean{PhyXMiniProblems.ProblemPhyXMini0894.AnswerChoice.forceInNewtons}
  \uses{def:physics:phyx-mini-0894:phyxminiproblems-problemphyxmini0894-answerchoice}
  Literal force magnitude in newtons printed beside an answer choice
\end{definition}
\begin{proof}
  This is the declared model definition of force In Newtons.
\end{proof}
\begin{definition}[recorded Dataset Answer]
  \label{def:physics:phyx-mini-0894:phyxminiproblems-problemphyxmini0894-recordeddatasetanswer}
  \lean{PhyXMiniProblems.ProblemPhyXMini0894.recordedDatasetAnswer}
  \uses{def:physics:phyx-mini-0894:phyxminiproblems-problemphyxmini0894-answerchoice}
  Answer label recorded by the supplied dataset metadata
\end{definition}
\begin{proof}
  This is the declared model definition of recorded Dataset Answer.
\end{proof}
\begin{lemma}[net Force Vector On Top In Newtons]
  \label{lem:physics:phyx-mini-0894:phyxminiproblems-problemphyxmini0894-netforcevectorontopinnewtons}
  \lean{PhyXMiniProblems.ProblemPhyXMini0894.netForceVectorOnTopInNewtons}
  \uses{def:physics:phyx-mini-0894:phyxminiproblems-problemphyxmini0894-forcevectorinnewtons, def:physics:phyx-mini-0894:phyxminiproblems-problemphyxmini0894-equilateralpointchargesetup, def:physics:phyx-mini-0894:phyxminiproblems-problemphyxmini0894-matchessuppliedequilateralchargefigure, def:physics:phyx-mini-0894:phyxminiproblems-problemphyxmini0894-hasequilateraltrianglegeometry, def:physics:phyx-mini-0894:phyxminiproblems-problemphyxmini0894-hasphysicalpointchargeparameters, def:physics:phyx-mini-0894:phyxminiproblems-problemphyxmini0894-usesroundedcoulombconstant, def:physics:phyx-mini-0894:phyxminiproblems-problemphyxmini0894-satisfiespointchargecoulombforcelaw, def:physics:phyx-mini-0894:phyxminiproblems-problemphyxmini0894-satisfieselectricforcesuperposition}
  The equal-magnitude pair forces have opposite vertical components and equal rightward horizontal components. Coulomb's law and superposition therefore give the following derived force vector on the top charge
\end{lemma}
\begin{proof}
  The conclusion follows from the governing relations and figure readouts named in the statement of net Force Vector On Top In Newtons.
\end{proof}
% --- Archon declaration topology end ---
% --- Archon physics formalization source end ---
